\documentclass[11pt,a4paper]{article}

% $Id$

% For up-to-date definitions of runlimits, etc.
% Configuration file for LaTeX
% $Id: config.template.tex 751 2005-05-21 13:44:43Z kink $

%%% AUTOGENERATE HEADER START %%%
%
% This configuration file was automatically generated
% with command './generate_config.sh tex'
% on Sun Sep 14 17:34:46 CEST 2008 on host 'delta'.
%
% Do not edit this file by hand! Instead, edit parts of this
% file which are outside the 'AUTOGENERATE HEADER' and
% 'GLOBAL CONFIG INCLUDE' tags in the templates 'config.template.*'.
%
% Configuration options inside 'GLOBAL CONFIG INCLUDE' tags
% should be edited in the main configuration file 'global.cfg'
% and then be included here by running 'make config' in the root
% of the system directory.
%
%%% AUTOGENERATE HEADER END %%%

%%% GLOBAL CONFIG INCLUDE START %%%
\def\DOMJUDGE_VERSION{3.0.0SVN}
\def\ROOT_BASE{/home/domjudge}
\def\SYSTEM_ROOT{/home/domjudge/system}
\def\OUTPUT_ROOT{/home/domjudge/output}
\def\INCOMINGDIR{/home/domjudge/output/incoming}
\def\SUBMITDIR{/home/domjudge/output/sources}
\def\JUDGEDIR{/home/domjudge/output/judging}
\def\LOGDIR{/home/domjudge/output/log}
\def\CHROOT_PREFIX{/home/domjudge/output/judging}
\def\WEBSERVER{domjudge.example.com}
\def\WEBBASEURI{http://domjudge.example.com/}
\def\SUBMITSERVER{domjudge.example.com}
\def\SUBMITPORT{9147}
\def\DBSERVER{domjudge.example.com}
\def\DBNAME{domjudge}
\def\SCP_HOST{localhost}
\def\RUNUSER{domjudge-run}
\def\VALID_USERS{domjudge-run}
\def\USE_CHROOT{1}
\def\CHROOT_SCRIPT{}
\def\COMPILETIME{30}
\def\SOURCESIZE{256}
\def\MEMLIMIT{524288}
\def\FILELIMIT{4096}
\def\PROCLIMIT{15}
\def\ENABLE_CMDSUBMIT_SERVER{1}
\def\ENABLE_WEBSUBMIT_SERVER{1}
\def\SUBMITCLIENT_METHOD{1}
\def\USERSUBMITDIR{.submit}
\def\USERPERMFILE{0600}
\def\USERPERMDIR{0700}
\def\WARN_MTIME{5}
\def\STRICTIPCHECK{0}
\def\VERIFICATION_REQUIRED{0}
\def\SHOW_COMPILE{2}
\def\SHOW_AFFILIATIONS{1}
\def\PENALTY_TIME{20}
\def\LANG_EXTS{C,c C++,cpp,cc,c++ Java,java Pascal,pas,p Haskell,hs,lhs Perl,pl Bash,sh}
\def\BEEP_CMD{/home/domjudge/system/bin/beep}
\def\BEEP_ERROR{-l 1000 -d 500 -f 800 -n -l 1000 -d 500 -f 1000 -n -l 1000 -d 500 -f 800 -n -l 1000 -d 500 -f 1000 -n -l 1000 -d 500 -f 800 -n -l 1000 -d 500 -f 1000 -n -l 1000 -d 500 -f 800 -n -l 1000 -d 500 -f 1000 -n -l 1000 -d 500 -f 800 -n -l 1000 -d 500 -f 1000}
\def\BEEP_WARNING{-l 1000 -d 500 -f 300 -n -l 1000 -d 500 -f 200 -n -l 1000 -d 500 -f 300 -n -l 1000 -d 500 -f 200}
\def\BEEP_SUBMIT{-f 400 -l 100 -n -f 400 -l 70}
\def\BEEP_ACCEPT{-f 400 -l 100 -n -f 500 -l 70}
\def\BEEP_REJECT{-f 400 -l 100 -n -f 350 -l 70}
\def\BALLOON_CMD{}
\def\SYSLOG{LOG_LOCAL0}
\def\DEBUG_PHP_NOTICE{1}
\def\DEBUG_TIMINGS{2}
\def\DEBUG_SQL{4}
\def\DEBUG_JUDGE{8}
\def\DEBUG{0}
%%% GLOBAL CONFIG INCLUDE END %%%


\usepackage[dutch]{babel}
\usepackage{a4wide}
\usepackage{verbatim}
\usepackage{moreverb}
\usepackage{amsmath,amssymb,amsthm}
\usepackage{epsfig}
\usepackage{bbm}
\usepackage{expdlist}
\usepackage{svn}

\setlength{\parindent}{0pt}
\setlength{\parskip}{1ex plus 0.5ex minus 0.2ex}

\newcommand{\DOMjudge}{\textsc{DOM}judge }

% We define our own \url command, because the one in the url package
% doesn't expand defines (present in systemconfig.tex).
\newcommand{\url}[1]{\texttt{#1}}


\title{Teamhandleiding \DOMjudge\\(versie \DOMJUDGE_VERSION)}

% For inclusion of the correct date
\SVNdate $Date$


\begin{document}

\begin{titlepage}
\maketitle
\vspace{1cm}
\begin{figure}[htbp]
\centerline{\includegraphics[height=7cm]{logos/DOMjudgelogo.eps}}
\end{figure}
\end{titlepage}

\newpage
\tableofcontents

\newpage
\section{Inleiding}

Deze documentatie beschrijft de onderdelen van het \DOMjudge
jurysysteem, die aan de teams ter beschikking staan om oplossingen in
te sturen, uitslagen te bekijken en ``clarification requests'' te
sturen.

Het systeem is voor een groot deel gebaseerd op een webinterface. Het
insturen van oplossingen gaat echter via een command-line programma.

\section{Samenvatting}

Hieronder staat kort samengevat de belangrijkste informatie. Dit is
bedoeld om snel aan de slag te kunnen. We raden echter ten zeerste
aan, dat minstens iemand binnen je team de complete handleiding
doorneemt, omdat daarin specifieke details van het jurysysteem staan,
die ook van belang kunnen zijn op het moment dat niet alles perfect
gaat (zoals een oplossing die fout is). \textbf{WEES GEWAARSCHUWD!!}

\subsection{Inlezen en wegschrijven}

Oplossingen moeten invoer en uitvoer lezen van `standard in'
(toetsenbord) en wegschrijven naar de `standard out' (beeldscherm).
Je hoeft dus nooit een bestand te openen. Zie bijlage
\ref{codeexamples} voor een aantal voorbeelden hiervan.

\subsection{Insturen van oplossingen}

Insturen van oplossingen gaat via het command-line programma
\texttt{submit}. Gebruik \texttt{submit <probleem>.<extensie>}, met
\texttt{<probleem>} het label van het probleem en \texttt{<extensie>}
een standaard extensie van je taal. Voor de complete documentatie en
alle opties, zie \texttt{submit --help}.

\subsection{Bekijken van scores, inzendingen, e.d.}

Het bekijken van inzendingen, scores en sturen en lezen van
``clarification requests'' gaat via een webinterface. Het adres is
\url{http://\WEBSERVER} en de pagina wijst zich verder vanzelf.

\newpage
\section{Oplossingen insturen: \texttt{submit}}\label{submit}

Het insturen van oplossingen voor problemen gebeurt via een
command-line interface: het programma \texttt{submit} (onder Windows
\texttt{submit.exe}).

\textbf{Syntax:} \texttt{submit [opties] bestandsnaam.ext}

Het submit programma haalt de naam (label) van het probleem uit
\texttt{bestandsnaam} en de de programmeertaal uit de extensie
\texttt{ext}. Dit kan handmatig aangepast worden met de opties
\texttt{-p probleemnaam} en \texttt{-l taalextensie}. Zie
\texttt{submit --help} voor een compleet overzicht van mogelijke
opties en extensies en een aantal voorbeelden. Let op, dat deze
help-tekst redelijk lang is en misschien niet op \'e\'en scherm
past. Gebruik dan bijvoorbeeld \texttt{submit --help | more} om alles
te lezen.

\texttt{submit} zal het bestand controleren op een aantal eigenschappen
en eventueel waarschuwingen geven, zoals wanneer het bestand al lange
tijd niet veranderd is of groter is dan de maximale source-code grootte.

Hierna geeft \texttt{submit} een kort overzicht met de details van de
inzending en vraagt om bevestiging. Controleer vooral of je het goede
bestand, probleem en taal hebt en druk dan op `y' om de oplossing in
te sturen. Als alles goed gaat, zal \texttt{submit} een melding geven
dat de inzending succesvol is. Indien niet, zal er een foutmelding
verschijnen.

Het submit programma maakt gebruik van een directory \texttt{\USERSUBMITDIR}
in de homedirectory van het account. Hier slaat het tijdelijk
bestanden op voor inzending en staat ook een logfile \texttt{submit.log}.
Verwijder deze directory niet of pas hem niet aan, omdat anders het
submit programma eventueel niet meer correct functioneert. Verder
is een ``public ssh-key'' van de jury in de ssh configuratie
toegevoegd. Deze is ook nodig voor het functioneren van \texttt{submit}.
 
\section{De uitslag bekijken van inzendingen}

Op de team webpagina staat een overzicht van je inzendingen.
Dit overzicht bevat alle relevante gegevens: de tijd van inzending, de
programmeertaal, het probleem en de status. Het adres van je
teampagina is \url{http://\WEBSERVER}. Hiervandaan kun je ook naar
de publieke pagina, waarop de scores van alle teams vermeld staan.

\subsection{Mogelijke uitslagen}

Op een ingestuurde oplossing zijn een aantal verschillende uitslagen
mogelijk. Hier volgen ze, met een korte beschrijving:

\begin{description}[\setleftmargin{4.5cm}]
\item[CORRECT]
Je oplossing was goed en je hebt dit probleem opgelost!

\item[COMPILER-ERROR]
Het compilen van je programma gaf een fout. Zie de inzendingsdetails
voor de precieze foutmelding. Compiler\emph{waarschuwingen} worden
niet als fouten gezien, maar kun je eventueel wel bekijken bij de
details van je inzending.

\item[TIMELIMIT]
Je programma draaide langer dan de maximaal toegestane tijd en is
afgebroken. Dit kan betekenen dat je programma ergens in een loop
blijft hangen, of dat je oplossing niet effici\"ent genoeg is.

\item[RUN-ERROR]
Je programma gaf een fout tijdens het uitvoeren. Dit kan verschillende
oorzaken hebben, zoals deling door nul, incorrecte geheugen adressering
(segfault, bijvoorbeeld door arrays buiten bereik te indiceren), te
veel geheugengebruik, enzovoort\dots\ 
Let ook op, dat je programma met een exitcode 0 eindigt!

\item[NO-OUTPUT]
Je programma gaf geen uitvoer. Let op dat je uitvoer naar standard
output schrijft!

\item[WRONG-ANSWER]
De uitvoer van je programma was niet correct. Het kan zijn dat je
oplossing niet correct is, maar let ook goed op dat je de antwoorden
precies zoals beschreven uitvoert: de uitvoer moet exact kloppen met
de antwoorden van de jury!
\end{description}

\section{Clarifications}

Communicatie met de jury loopt door middel van clarifications, deze
komen op je teampagina te staan.  Boven aan de pagina de gedane
clarifications, daar onder je requests.

Je kan vragen aan de jury stellen door middel van het doen van een
``Clarification Request'', de link hiervoor bevindt zich onderaan de
clarifications-pagina.  Je vraag zal alleen bij de jury aan komen, zij
zullen deze zo snel mogelijk en adequaat beantwoorden. Antwoorden die
voor iedereen relevant zijn zullen naar iedereen gestuurd worden.

Als je een clarification ontvangt van de jury, dan wordt dat
automatisch gemeld door middel van een popup scherm. Hiervoor moet je
wel de teampagina open hebben staan in een webbrowser, die geen popup
schermen blokkeert.

\section{Hoe worden opgaven beoordeeld?}

Het \DOMjudge jurysysteem is volledig geautomatiseerd. Dit betekent
dat er (in principe) geen menselijke interactie is tijdens de
beoordeling. Het beoordelen gebeurt via de volgende stappen:

\subsection{Insturen}

Via het \texttt{submit} (zie \ref{submit}) programma kun je een
oplossing voor een opgave insturen. Let op dat je de source-code van
je programma moet insturen (en dus niet een gecompiled programma of de
uitvoer van je programma). \texttt{submit} maakt dan een verbinding
met een server en die kopieert je programma dan naar de
jury-computers.

Dan komt je programma in de wachtrij te staan, om gecompiled,
uitgevoerd en de uitvoer getest te worden op \'e\'en van de
jury-computers.

\subsection{Compileren}

Je programma wordt op een jury-computer onder Linux gecompileerd.
Als je een andere compiler of besturingssysteem gebruikt dan de jury,
moet dat in principe geen probleem zijn, maar let wel op dat
je geen compiler/systeem-specifieke dingen gebruikt (bij een
compileerfout kun je altijd de foutmelding bekijken).

Verder wordt door in het jurysysteem alles statisch gecompileerd, maar
wat betreft foutmeldingen e.d. zou dit niet relevant moeten zijn.

\subsection{Testen}

Als je programma succesvol gecompileerd is, wordt het gedraaid en de
uitvoer vergeleken met de correcte uitvoer van de jury. Er wordt eerst
gecontroleerd of je programma correct ge\"eindigd is: als je programma
met een fout eindigt en het goede antwoord geeft, krijg je toch een
run-error! Er zijn een aantal beperkingen die aan je programma
opgelegd worden. Als je programma die overschrijdt, wordt het ook
afgebroken met een fout, zie (\ref{runlimits}).

Verder moet de uitvoer van jouw programma exact overeenkomen met de
uitvoer van de jury. Let dus goed op, dat je de uitvoerspecificatie
volgt. In gevallen waarin er niet \'e\'en unieke uitvoer is (zoals bij
floating point antwoorden) kan de jury een aangepaste beoordeling
hiervoor maken.

\subsection{Beperkingen}\label{runlimits}

Om hacken tegen te gaan, het jurysysteem stabiel te houden en iedereen
duidelijke, gelijke omstandigheden te geven, zijn er een aantal
beperkingen die aan iedere ingestuurde oplossing opgelegd worden:

\begin{description}[\setlabelphantom{aantal processen}]
\item[compile-tijd]
Je programma mag er maximaal \COMPILETIME\ seconden over doen om te
compileren. Daarna wordt het compileren afgebroken en levert dit een
compileerfout op. Dit zou in de praktijk nooit een probleem mogen
opleveren. Mocht dit toch gebeuren bij een normaal programma, laat het
dan de organisatie weten.

\item[source grootte]
De sourcecode van je programma mag maximaal \SOURCESIZE\ kilobytes
groot zijn. Bij het inzenden zul je hier al voor gewaarschuwd worden,
bij compileren levert dit ook een compileerfout op.

\item[geheugen]
Je programma heeft tijdens het draaien maximaal \MEMLIMIT\ kilobytes
geheugen ter beschikking. Let op dat dit totaal geheugen is (inclusief
programmacode, statisch en dynamisch gedefini\"eerde variabelen,
stack, \dots)! Als je programma meer probeert te gebruiken, zal het
waarschijnlijk afbreken, zodat dit een ``RUN-ERROR'' geeft.

\item[bestandsgrootte]
Hoewel je geen bestanden mag schrijven, wordt de standaard (fout)uitvoer
(stout/stderr) wel naar een bestand geschreven. De maximale grootte
van deze bestanden is \FILELIMIT\ kilobytes (per bestand).

\item[aantal processen]
Het is niet de bedoeling dat je programma meerdere processen (threads)
start. Dit heeft ook geen zin, want je programma heeft precies 1
processor volledig tot zijn beschikking. Om de stabiliteit van het
jurysysteem te bevorderen, kun je maximaal \PROCLIMIT\ processen
tegelijk draaien (inclusief de processen waardoor je programma
gestart is).

Mensen die nooit met meerdere processen geprogrammeerd hebben (of
niet weten wat dat is), hoeven zich geen zorgen te maken: standaard
draait een gecompileerd programma in \'e\'en proces.

\end{description}


\newpage
\appendix

\section{Code voorbeelden}\label{codeexamples}

Hieronder staan een aantal voorbeelden van code om de invoer van een
probleem in te lezen en de uitvoer weg te schrijven. De manier van
in- en uitvoer lezen en wegschrijven in deze opgave, is zoals dat bij
veel programmeerwedstrijden gebeurt (o.a. de ACM regionals), maar kan
bij jullie wedstrijd anders zijn.

De code hoort bij de volgende probleembeschrijving: gegeven een string
met een naam (\'e\'en woord), voer de string ``Hello $<$naam$>$!'' uit.

Dit probleem zou de volgende in- en uitvoer kunnen hebben:

\begin{tabular}{|p{0.5\textwidth}|p{0.5\textwidth}|}
\hline
\textbf{Invoer} & \textbf{Uitvoer} \\
\hline
\verbatiminput{example.in} &
\verbatiminput{example.out} \\
\hline
\end{tabular}

Let op dat het op de eerste regel het getal 3 aangeeft, dat er 3
testgevallen volgen.

Een oplossing voor dit probleem in C:
\listinginput{1}{example.c}

Let op de \texttt{return 0;} aan het einde, zodat we geen
run-error krijgen!

\newpage

Een oplossing in C++ kan als volgt:
\listinginput{1}{example.cc}

Een oplossing in Pascal:
\listinginput{1}{example.pas}

\newpage

En tenslotte een oplossing in Java:
\listinginput{1}{example.java}

\end{document}
